\documentclass[14pt,a1paper,landscape]{tikzposter}

%% Tikzposter is highly customizable: please see
%% https://bitbucket.org/surmann/tikzposter/downloads/styleguide.pdf

%% Available themes: see also
%% https://bitbucket.org/surmann/tikzposter/downloads/themes.pdf
% \usetheme{Default}
% \usetheme{Rays}
% \usetheme{Basic}
% \usetheme{Simple}
 \usetheme{Envelope}
% \usetheme{Wave}
% \usetheme{Board}
% \usetheme{Autumn}
% \usetheme{Desert}

%% Further changes to the title etc is possible
% \usetitlestyle{Default}
% \usetitlestyle{Basic}
 \usetitlestyle{Empty}
% \usetitlestyle{Filled}
% \usetitlestyle{Envelope}
% \usetitlestyle{Wave}
% \usetitlestyle{verticalShading}


%\usepackage{fontspec}
%\setmainfont{FreeSerif}
%\setsansfont{FreeSans}

\usepackage [sfdefault]{cabin}
\usepackage[T1]{fontenc}
\usepackage{natbib}

\author{Eric Rexstad, David L. Miller, Laura Marshall and Len Thomas}
\title{\parbox{0.6\linewidth}{\centering Migrating distance sampling projects from Distance for Windows to the Distance R package}}
\institute{Centre for Research into Ecological and Environmental Modelling University of St Andrews}


% borrowed from https://tex.stackexchange.com/questions/185166/tikzposter-4-logos-title-2-at-each-side/185278#185278

%\makeatletter
%\def\TP@titlegraphictotitledistance{-4cm}
%\settitle{ \centering \vbox{
%		\@titlegraphic \\ [\TP@titlegraphictotitledistance] 
%		\centering
%		\color{titlefgcolor} {\bfseries \Huge \sc \@title \par}
%		\vspace*{1em}
%		{\huge \@author \par} \vspace*{1em} {\LARGE \@institute}
%}}
%\makeatother

%\setlength{\columnsep}{2cm}
%
%\title{\parbox{0.5\linewidth}{\centering My  really  really  really  really  really  really  really  really  long poster TITLE.}}
%\author{Name}
%\titlegraphic{
%	\includegraphics[width=8cm, height=4cm]{ISEC_Logo.png}
%	\hfill 
%%	\includegraphics[width=9cm,height=8cm]{distancelogo.png}
%	\\% \vspace{5pt} \\
%	\includegraphics[width=3cm, height=6cm]{02-foundation-vertical-white.png}
%	\hfill 
%	\includegraphics[width=10.5cm,height=3cm]{CREEM_Logo_trans.png}
%}


%% Original, functional way to place single logo in middle
%\titlegraphic{\includegraphics[width=7cm]{ISEC_Logo.png}}
%% Uncomment to switch off tikzposter footer
\tikzposterlatexaffectionproofoff

\begin{document}
\maketitle

\begin{columns}

\column{0.333}

\block{Introduction}{
The Distance software \citep{thomas_distance_2010} has been downloaded >40,000 times in its 20-year history.  Much of the underlying machinery is written in R.  For some users, there may be benefits to performing the analysis with the underlying R code, rather than working with the graphical user interface (GUI).
	\begin{tikzfigure}[Traditional Distance for Windows interface (left) and distance sampling analysis in R (right).]
%		\includegraphics[width=0.45\columnwidth]{D7-analysis-grab.PNG}
		\includegraphics[width=0.95\linewidth]{merged-grabs.png}
%		\includegraphics[width=0.45\columnwidth]{R-analysis-grab.PNG}
	\end{tikzfigure}

Challenges hindering the transition between analysis with the GUI and analyses in R are two-fold:
\begin{itemize}
	\item Legacy data reside in Distance (GUI) projects, unavailable for importing into R, and
	\item Analyses that are easily described using the GUI may be difficult to specify, particularly if analyst is not proficient in R.
\end{itemize}
}

\note[rotate=8, connection, width = 7cm, % targetoffsety=-2cm,
% roundedcorners=15, targetoffsetx=2cm
]{How to make this transition?}

\block{How to bridge between the two?}{

\begin{tikzfigure}[Database structure of a Distance project.]
	\includegraphics[width=0.3\linewidth]{dst-file-structure.PNG}
\end{tikzfigure}

Distance project \textbf{\texttt{(.dst)}} file contains:

\begin{enumerate}
\item Data
\item model definitions
\item analysis results
\end{enumerate}

within an Access database.

\begin{itemize}
	\item Trick is to extract contents of (2) and translate into R code
	\item Target the results of previous step to contents of (1) the data
	\item Perhaps contrast results of R analysis with results stored in the Access database
\end{itemize}


}

\column{0.3333}
\block{Applications--Legacy projects}{
\begin{itemize}
\item We try to do this. what?
\item we try to do that.
\end{itemize}
}

\block{Learning structure of R interface}{%
\begin{tikzfigure}[The organisation for which we work]
\includegraphics[width=0.35\linewidth]{CREEM_Logo.jpg}
\end{tikzfigure}
}

\column{0.3333}
\block{Comparative analysis of difficult data}{
\begin{itemize}
\item We try to do this. what?
\item we try to do that.
\end{itemize}
}

\block{Caveats}{
	\texttt{readdst} is not able to translate all GUI analyses into R code.  Current limitations are inability to translate
	\begin{itemize}
		\item analyses using the \texttt{dsm}, \texttt{mads} and \texttt{Dssim} engines,
		\item analyses using post-stratification and
		\item bootstraps for variance estimation.
	\end{itemize}
}

\block{Request for additional datasets}{
\begin{itemize}
\item We try to do this. what?
\item we try to do that.
\item We try to do this. what?
\item we try to do that.
\item We try to do this. what?
\item we try to do that.
\end{itemize}
}

\block{Additional information}{

\innerblock{Literature}{
\nocite{*} % Insert publications even if they are not cited in the poster
\small{\bibliographystyle{humannat}
	\bibliography{isecposter}\vspace{0.75in}}
}
\item QR codes to package/website/bioRxiv

}

\end{columns}


\end{document}